EJERCICIO: Juego de Memoria con Emojis (POO)

Desarrollar una aplicación web utilizando HTML, CSS y JavaScript puro que implemente un juego
de memoria.

Requisitos:
- Tablero de 10x10 (100 cartas).
- 50 pares de emojis (cada uno repetido 2 veces).
- Cartas ocultas inicialmente.
- Mostrar emoji al hacer click.

Lógica:
- Solo 2 cartas por turno.
- Si coinciden: quedan visibles.
- Si no coinciden: se ocultan luego de 1 segundo.
- No permitir seleccionar cartas ya resueltas.

Fin del juego:
- Cuando todas las cartas estén descubiertas mostrar: 'Juego terminado'.
- Y tambien se tien que mostrar la cantidad de intentos que isite.

Programación orientada a objetos:
- Clase Carta: emoji, estado
- Clase Juego: lógica, mezcla, control de turnos

Restricciones:
- No usar librerías externas.
- No usar alert ni prompt.
- Usar solo DOM.